\documentclass[10pt,a4paper]{amsart}
\usepackage[margin=19mm]{geometry}
\usepackage{amsmath,amssymb,mathtools}
\usepackage[scr=rsfs,cal=euler]{mathalfa}
\usepackage{xfrac}
\usepackage[unicode,hidelinks,bookmarks=false]{hyperref}
\newcommand{\ie}{i.e.\ }
\newcommand{\cf}{cf.\ }
\newcommand{\resp}{respectively\ }
\newcommand{\Subsection}{\subsection}
\providecommand{\nobreakdash}{\nobreak\hbox{-}}
\setlength{\emergencystretch}{2em}
\multlinegap0pt
\usepackage{bm}
\usepackage{bbm}
\usepackage{dsfont}
\usepackage{xspace}
\usepackage{cancel}
\usepackage{mathtools}
\usepackage{amsthm}
\makeatletter
\@ifundefined{lemma}{\newtheorem{lemma}{Lemma}}{}
\makeatother
\title[Non-convergence and convergence of bounded solutions to semi]{Non-convergence and convergence of bounded solutions to semilinear wave equation with dissipative boundary condition}
\author{Zhe Jiao, Xiao Li}
\date{Source: arXiv:2506.13165v1 (CC BY 4.0)}
\begin{document}
\maketitle

\noindent\emph{Source and licence.} Non-convergence and convergence of bounded solutions to semilinear wave equation with dissipative boundary condition. Version arXiv:2506.13165v1 (CC BY 4.0), distributed under the Creative Commons Attribution 4.0 International licence, as recorded in the arXiv licence metadata for this exact version and shown on the abstract page https://arxiv.org/abs/2506.13165v1. Full rights notice: licenses/arxiv-2506.13165-CC-BY-4.0.txt.\par\smallskip

\noindent\textit{Editorial scope.} Counted unit: the Lemma whose source label is \texttt{decomposition\_1} in the article (source0.tex lines 997--1007, source label \texttt{decomposition\_1}), delivered together with the article's own complete original proof of it (source0.tex lines 1008--1142, one \texttt{proof} environment of 135 lines). No mathematics is written, repaired or re-derived: statement and proof are the article's own text line for line. Required context retained uncounted: cmr\_ode (lines 902--904); cmr\_ode1 (lines 906--915); cmr\_ode2 (lines 917--922); cmr\_estimate (lines 924--926). Every definition, display and label of those lines is retained unchanged. Omitted from this package and not redistributed: 1--901, 905--905, 916--916, 923--923, 927--996, 1143--1923. Nothing inside a retained excerpt is omitted or altered: every retained body is delivered exactly as the article's lines are written, so no omission receipt of any kind accompanies this package. Citations: the retained excerpts carry no citation command at all, so no bibliography entry of the article is transcribed and no reference is redistributed. Reference adaptation: none was needed and none was made; every reference inside the delivered unit resolves inside it, so no unresolved reference or citation is produced. Printed slips kept literal: no printed word, symbol, hypothesis or number of the counted statement or of its proof was altered. Layout and class: the article's own class is \texttt{elsarticle} with packages including geometry, amssymb, amsfonts, verbatim, xcolor, amsmath, amsthm, graphicx, subfigure, booktabs, adjustbox, multirow, xy, pifont; the wrapper is the lane's pinned \texttt{amsart} editorial wrapper at 10pt on A4, which restates the theorem-like environments the retained text needs and adds the article's own macro definitions that the retained text needs (none), with extra wrapper packages (bm, bbm, dsfont, xspace, cancel, mathtools). The delivered numbering is the unit's own. Assets: no retained span contains a figure environment, a graphics-inclusion command, a PDF-page inclusion, a table or a TikZ picture; no image file is copied into this package, the package declares no asset dependency and no image file is redistributed with it. Licence: version arXiv:2506.13165v1 of this article is distributed under the Creative Commons Attribution 4.0 International licence, as recorded in the arXiv licence metadata for this exact version and shown on the abstract page https://arxiv.org/abs/2506.13165v1. A full rights notice is delivered at licenses/arxiv-2506.13165-CC-BY-4.0.txt. Characters: every retained excerpt is pure ASCII, so the delivered engine, XeLaTeX with its default Latin Modern text font, reports no missing character.
\begin{equation} \label{cmr_ode}
\dot{\xi}(t) = h(\xi(t)), \quad \xi(0) = \xi_0
\end{equation}

\begin{equation} \label{cmr_ode1}
\left(
\begin{array}{c}
h(\xi)\cdot\psi \\
0
\end{array}
\right)
=
-\mathcal{P}\mathcal{F}(U).
\end{equation}

\begin{equation}\label{cmr_ode2}
\begin{split}
	\nabla_{\xi} U \cdot h(\xi) + \tilde{\mathcal{A}}U(t) + \mathcal{F}(U(t)) &= 0,\\
	\nabla_{\xi} \Lambda \cdot h(\xi) + \tilde{\mathcal{A}}\Lambda(\xi) + (\mathcal{I}-\mathcal{P})\mathcal{F}(U(\xi)) &= 0.
\end{split}
\end{equation}

\begin{equation}\label{cmr_estimate}
	\|\Lambda\|_{C^{j-1}_{\mathrm{b}}(\mathbb{R}^2\rightarrow \mathcal{D})} \leqslant C\|\mathcal{F}\|_{C^{1}_{\mathrm{b}}(\mathcal{D}\rightarrow \mathcal{D})}
\end{equation}

\begin{lemma}\label{decomposition_1}
Any equilibrium of~\eqref{cmr_ode} is a critical point of the functional $\Xi$. Moreover, if $f_0$ is of class $ C^3_{\mathrm{b}}(\bar{\Omega}\times\mathbb{R})$ and $\|f_a\|_{C^2_{\mathrm{b}}(\bar{\Omega}\times\mathbb{R})}< c^{\ast} < c_3$ with $c^{\ast}$ sufficiently small, then one has
\[
	\Xi(\xi) = \Upsilon_1(\xi) +  \int_{\Omega}F_a(x, \xi\cdot\psi)\mathrm{d}x, \quad \xi\in\mathbb{R}^2,
\]
where $\Upsilon_1(\xi)$ is a $C^2$ function on $\mathbb{R}^2$ satisfying
\begin{equation} \label{estimates_lyapunov_3}
|\Upsilon_1(\xi)| + |\partial_{\xi_i}\Upsilon_1(\xi)|\leqslant C\max\{1, \|f_a\|^2_{C^3_{\mathrm{b}}(\bar{\Omega}\times\mathbb{R})}\}\left( \sup_{x\in\Omega}|f_a(x, \xi\cdot\psi)|^2 + \sup_{x\in\Omega}|f^{\prime}_a(x, \xi\cdot\psi)| \right)
\end{equation}
with $i=1, 2$ and $C$ a positive constant.
\end{lemma}

\begin{proof}
To facilitate reading, we divide the proof into the following three steps.

\textbf{Step 1}. Combining integration by parts and the fact that $\xi\cdot\psi\in\mathrm{Ker}B$ gives
\begin{equation} \label{lyapunov_1}
\begin{split}
\Xi(\xi) &= \frac{1}{2}\left\|\left(\xi\cdot\psi + \Lambda_1(\xi), \Lambda_2(\xi)\right)^{\top}\right\|^2_{\mathcal{H}} + \int_{\Omega}F_0(x, \xi\cdot\psi + \Lambda_1(\xi))\mathrm{d}x\\
&= \frac{1}{2}\left(\int_{\Omega}|\nabla (\xi\cdot\psi + \Lambda_1(\xi))|^2\mathrm{d}x + \int_{\Gamma}|\xi\cdot\psi +\Lambda_1(\xi)|^2\mathrm{d}S + \int_{\Omega}|\Lambda_2(\xi)|^2\mathrm{d}x\right) \\
&\quad + \int_{\Omega}F_0(x, \xi\cdot\psi + \Lambda_1(\xi))\mathrm{d}x\\
&= \frac{1}{2}\left(\int_{\Omega}|\nabla  \Lambda_1(\xi)|^2\mathrm{d}x - a(x)\int_{\Omega}|  \Lambda_1(\xi)|^2\mathrm{d}x + \int_{\Gamma}|\Lambda_1(\xi)|^2\mathrm{d}S + \int_{\Omega}|\Lambda_2(\xi)|^2\mathrm{d}x\right) \\
&\quad+\int_{\Omega}F_a(x, \xi\cdot\psi + \Lambda_1(\xi))\mathrm{d}x
\end{split}
\end{equation}
with $F_a(x, u)=\int_{0}^{u}f_a(x, s)\mathrm{d}s$.
Differentiating~\eqref{lyapunov_1} with respect to $\xi_i$, $i=1, 2$, we obtain
\begin{equation}\label{lyapunov_2}
\begin{split}
\partial_{\xi_i}\Xi(\xi) &=  -\int_{\Omega}\mathcal{B}\Lambda_1(\xi)\partial_{\xi_i}\Lambda_1(\xi)\mathrm{d}x + \int_{\Omega}\Lambda_2(\xi)\partial_{\xi_i}\Lambda_2(\xi)\mathrm{d}x-\int_{\Gamma} \Lambda_2(\xi)\partial_{\xi_i}\Lambda_1(\xi)\mathrm{d}S\\
&\quad + \int_{\Omega}f_a(x, \xi\cdot\psi + \Lambda_1(\xi))(\psi_i + \partial_{\xi_i}\Lambda_1(\xi))\mathrm{d}x\\
&=  -\int_{\Omega}[\mathcal{B}\Lambda_1(\xi)-f_a(x, \xi\cdot\psi + \Lambda_1(\xi))](\psi_i + \partial_{\xi_i}\Lambda_1(\xi))\mathrm{d}x - \int_{\Gamma}\psi_i \Lambda_2(\xi)\mathrm{d}S\\
&\quad + \int_{\Omega}\Lambda_2(\xi)\partial_{\xi_i}\Lambda_2(\xi)\mathrm{d}x-\int_{\Gamma} \Lambda_2(\xi)\partial_{\xi_i}\Lambda_1(\xi)\mathrm{d}S.
\end{split}
\end{equation}
The first equality of~\eqref{cmr_ode2} can be written in component wise
\begin{equation}\label{cmr_ode3}
\begin{split}
	  h(\xi)\cdot \psi + \nabla_{\xi}\Lambda_1 \cdot h(\xi) &= \Lambda_2(\xi),\\
	\nabla_{\xi}\Lambda_2 \cdot h(\xi) &= \mathcal{B}\Lambda_1 - T^\ast(T\Lambda_2) - f_a(x, h(\xi)\cdot \psi +\Lambda_1).
\end{split}
\end{equation}
For every $\xi_0$, $h(\xi_0) = 0$ implies from~\eqref{cmr_ode3} that 
\[
\Lambda_2(\xi_0) = 0, \quad \mathcal{B}\Lambda_1(\xi_0) - f_a(x, \Lambda_1(\xi_0)) = 0.
\]
Consequently, it deduce from~\eqref{lyapunov_2} that $\nabla_{\xi}\Xi(\xi_0) = 0$. This completes the first part of the lemma.

\textbf{Step 2}. Due to the following equality
\begin{equation*}
\begin{split}
 \int_{\Omega}F_a(x, \xi\cdot\psi + \Lambda_1(\xi))\mathrm{d}x -\int_{\Omega}F_a(x, \xi\cdot\psi)\mathrm{d}x &=  \int_{\Omega}\int_{\xi\cdot\psi}^{\xi\cdot\psi + \Lambda_1(\xi)}f_a(x, y)\mathrm{d}y\mathrm{d}x\\
 &=\int_{\Omega}\Lambda_1(\xi))\int^{1}_{0}f_a(x, \xi\cdot\psi + s\Lambda_1(\xi))\mathrm{d}s\mathrm{d}x,	
\end{split}
\end{equation*}
then from~\eqref{lyapunov_1} the Lyapunov functional $\Xi(\xi)$ can be reformulated as
\begin{equation*}
\Xi(\xi) =  \Upsilon_1(\xi) +  \int_{\Omega}F_a(x, \xi\cdot\psi)\mathrm{d}x
\end{equation*}
with
\begin{equation}\label{lyapunov_3}
\begin{split}
\Upsilon_1(\xi) &:= \frac{1}{2}\left(\int_{\Omega}|\nabla  \Lambda_1(\xi)|^2\mathrm{d}x - a(x)\int_{\Omega}|  \Lambda_1(\xi)|^2\mathrm{d}x + \int_{\Gamma}|\Lambda_1(\xi)|^2\mathrm{d}S + \int_{\Omega}|\Lambda_2(\xi)|^2\mathrm{d}x\right) \\
&\qquad + \int_{\Omega}\Lambda_1(\xi))\int^{1}_{0}f_a(x, \xi\cdot\psi + s\Lambda_1(\xi))\mathrm{d}s\mathrm{d}x. 
\end{split}
\end{equation}
From differentiating~\eqref{lyapunov_3} with respect to $\xi_i$ it deduces that
\begin{equation}\label{diff_lyapunov_3}
\begin{split}
\partial_{\xi_i}\Upsilon_1(\xi) &=  -\int_{\Omega}\mathcal{B}\Lambda_1(\xi)\partial_{\xi_i}\Lambda_1(\xi)\mathrm{d}x + \int_{\Omega}\Lambda_2(\xi)\partial_{\xi_i}\Lambda_2(\xi)\mathrm{d}x-\int_{\Gamma} \Lambda_2(\xi)\partial_{\xi_i}\Lambda_1(\xi)\mathrm{d}S\\
&\quad + \int_{\Omega}\Lambda_1(\xi))\int^{1}_{0} f^{\prime}_a(x, \xi\cdot\psi + s\Lambda_1(\xi))(\psi_i + s\partial_{\xi_i}\Lambda_1(\xi))\mathrm{d}s\mathrm{d}x\\
&\quad + \int_{\Omega}\partial_{\xi_i}\Lambda_1(\xi))\int^{1}_{0}f_a(x, \xi\cdot\psi + s\Lambda_1(\xi))\mathrm{d}s\mathrm{d}x.
\end{split}
\end{equation}

To estimate $\Upsilon_1(\xi)$ and $\partial_{\xi_i}\Upsilon_1(\xi)$, we need the following estimates
\begin{equation}\label{estimates_lyp_3_1}
\begin{split}
	\| \Lambda_1(\xi)\|_{H^2} + \| \Lambda_2(\xi)\|&\leqslant C_1\|\tilde{f}_a(\xi\cdot\psi)\|,\\
	\|\tilde{f}_a(x, \xi\cdot\psi + s\Lambda_1(\xi))\|&\leqslant C_2 \|\tilde{f}_a(\xi\cdot\psi)\|,\\
	\|\partial_{\xi_i} \Lambda_1(\xi)\|_{H^2} + \| \partial_{\xi_i}\Lambda_2(\xi)\|&\leqslant C_3\left(\|f_a\|_{C^3_{\mathrm{b}}(\bar{\Omega}\times\mathbb{R})}\|\tilde{f}_a(\xi\cdot\psi)\| + \|\tilde{f}^{\prime}_a(\xi\cdot\psi)\psi_i \| \right), \\
	\|\partial_{\xi_i}\tilde{f}_a(x, \xi\cdot\psi + s\Lambda_1(\xi))\|&\leqslant C_4\left(\|f_a\|_{C^3_{\mathrm{b}}(\bar{\Omega}\times\mathbb{R})}\|\tilde{f}_a(\xi\cdot\psi)\| + \|\tilde{f}^{\prime}_a(\xi\cdot\psi)\psi_i \| \right).
\end{split}
\end{equation}
By~\eqref{lyapunov_3} and the first two estimates of~\eqref{estimates_lyp_3_1} we get
\[
|\Upsilon_1(\xi)| \leqslant C_5\|\tilde{f}_a(\xi\cdot\psi)\|^2.
\]
Similarly, by~\eqref{diff_lyapunov_3} and the last two estimates of~\eqref{estimates_lyp_3_1} we have
\[
|\partial_{\xi_i}\Upsilon_1(\xi)| \leqslant C_6\left(\|f_a\|^2_{C^3_{\mathrm{b}}(\bar{\Omega}\times\mathbb{R})}\|\tilde{f}_a(\xi\cdot\psi)\|^2 + \|\tilde{f}^{\prime}_a(\xi\cdot\psi)\psi_i \|^2 \right).
\]
We also note that
\begin{equation*}
	\|\tilde{f}_a(\xi\cdot\psi)\|\leqslant C_7\sup_{x\in\Omega}|f_a(x, \xi\cdot\psi)|,  \quad \|\tilde{f}^{\prime}_a(\xi\cdot\psi)\|\leqslant C_8 \sup_{x\in\Omega}|f^{\prime}_a(x, \xi\cdot\psi)|.
\end{equation*}
Thus, the expected inequality~\eqref{estimates_lyapunov_3} can be obtained.

\textbf{Step 3}. The remaining part of this proof is to show the estimates~\eqref{estimates_lyp_3_1} hold true. 

Let $\tilde{\mathcal{B}}$ be the restriction of the operator $\mathcal{B}$ to $(\mathrm{Ker}\mathcal{B})^{\bot}$. Then its inverse operator $\tilde{\mathcal{B}}^{-1}: (\mathrm{Ker}\mathcal{B})^{\bot}\rightarrow H^2(\Omega)$ is bounded. From~\eqref{cmr_ode3} we have
\begin{equation} \label{eq:lambdas}
\begin{split}
	\Lambda_1(\xi) &= \tilde{\mathcal{B}}^{-1}\left(\nabla_{\xi}\Lambda_2 \cdot h(\xi) + T^\ast(T\Lambda_2) + f_a(x, h(\xi)\cdot \psi + \Lambda_1)\right),\\
	\Lambda_2(\xi) &= h(\xi)\cdot \psi + \nabla_{\xi}\Lambda_1 \cdot h(\xi).
\end{split}
\end{equation}

By~\eqref{cmr_ode1} and ~\eqref{cmr_estimate}, the second equality of~\eqref{eq:lambdas} yields
\begin{equation} \label{eq:lambdas_2}
\begin{split}
\|\Lambda_2(\xi)\|&\leqslant \|h(\xi)\cdot \psi \| + |\nabla_{\xi}\Lambda_1|\|h(\xi)\|\\
&\leqslant C_8\|\tilde{f}_a(\xi\cdot\psi + \Lambda_1(\xi))\|.
\end{split}
\end{equation}
A similar estimate in the first equality of~\eqref{eq:lambdas} gives
\begin{equation} \label{eq:lambdas_1}
\begin{split}
\|\Lambda_1(\xi)\|_{H^2} &\leqslant C_9\left(\|\tilde{f}_a(\xi\cdot\psi + \Lambda_1(\xi))\| + \|\Lambda_2(\xi)\|\right)\\
	&\leqslant C_{10}\|\tilde{f}_a(\xi\cdot\psi + \Lambda_1(\xi))\|.
\end{split}
\end{equation}
If $\|f_a\|_{C^2_{\mathrm{b}}(\bar{\Omega}\times\mathbb{R})}< c^{\ast}$ with $c^\ast$ small enough, then by the mean-value theorem we have
\begin{equation} \label{estimate:f_a}
C_{10}\|\tilde{f}_a(\xi\cdot\psi + \Lambda_1(\xi))\|\leqslant C_{10}\|\tilde{f}_a(\xi\cdot\psi)\| +\frac{1}{2}\|\Lambda_1(\xi)\|.
\end{equation}
Thus, from~\eqref{eq:lambdas_1}, \eqref{eq:lambdas_2} and~\eqref{estimate:f_a} it implies the first two estimates of~\eqref{estimates_lyp_3_1}.

By differentiating \eqref{eq:lambdas} with respect to $\xi_i$, it implies from~\eqref{cmr_ode1} and~\eqref{cmr_estimate} that
\begin{equation} \label{estimate:partial_xi}
	\|\partial_{\xi_i} \Lambda_2(\xi)\|\leqslant C_{11}\mathcal{Q}, \quad \|\partial_{\xi_i} \Lambda_1(\xi)\|_{H^2} \leqslant C_{12}\left( \mathcal{Q} + \|\partial_{\xi_i} \Lambda_2(\xi)\|\right)
\end{equation}
with
\begin{equation*}
	\mathcal{Q}:= \|\tilde{f}^{\prime}_a(\xi\cdot\psi + \Lambda_1(\xi))(\psi_i + \partial_{\xi_i}\Lambda_1(\xi))\| +\|f_a\|_{C^3_{\mathrm{b}}(\bar{\Omega}\times\mathbb{R})}\|\tilde{f}_a(\xi\cdot\psi + \Lambda_1(\xi))\|.
\end{equation*}
Combining the first two estimates of~\eqref{estimates_lyp_3_1} and~\eqref{estimate:partial_xi} with the mean-value theorem gives 
\begin{equation*}
\begin{split}
	&\|\partial_{\xi_i} \Lambda_1(\xi)\|_{H^2} + \|\partial_{\xi_i} \Lambda_2(\xi)\| \\
	&\leqslant C_{13}\mathcal{Q} \\
	&\leqslant C_{14}\left(\|\tilde{f}^{\prime}_a(\xi\cdot\psi)\psi_i \| + \|f_a\|_{C^3_{\mathrm{b}}(\bar{\Omega}\times\mathbb{R})}\|\Lambda_1(\xi)\| + \|f_a\|_{C^3_{\mathrm{b}}(\bar{\Omega}\times\mathbb{R})}\|\tilde{f}_a(\xi\cdot\psi)\| + c^\ast\|\partial_{\xi_i} \Lambda_1(\xi)\|_{H^2} \right)\\
	&\leqslant C_{15}\left(\|\tilde{f}^{\prime}_a(\xi\cdot\psi)\psi_i \| + \|f_a\|_{C^3_{\mathrm{b}}(\bar{\Omega}\times\mathbb{R})}\|\tilde{f}_a(\xi\cdot\psi)\| + c^\ast\|\partial_{\xi_i} \Lambda_1(\xi)\|_{H^2} \right),
\end{split}
\end{equation*}
which implies the third estimate of~\eqref{estimates_lyp_3_1} by choosing $c^\ast$ small enough. The last estimate of~\eqref{estimates_lyp_3_1} follows from the third estimate and the mean-value theorem. This concludes the proof.
\end{proof}

\end{document}
